\chapter{SPARK simulation: construction, results and comparison}
\label{sec:towers-spark-study}

\section{Executed campaign and experimental coverage}

The native experiment is an Ada implementation with SPARK-analyzed counting,
numerical and driver code. Its retained run uses \(T_S=1{,}000{,}000\)
shared pairs, seed \texttt{20261030}, and the four original rules E2, E0.1,
Rare and Mixed. There is one pseudorandom stream: the four rows classify the
same pairs. Four million rule classifications therefore do not constitute
four million independent pairs. OR2 is calculated by the native numerical
code, but is not a fifth native Monte Carlo row; the Dask experiment supplies
that additional simulated rule.

The original 31 October output is the historical input pinned by
\path{reporting/run.json}. The later 4 November source-pinned campaign
retains the same four-row CSV, byte for byte, under both Ghost policies.
This report uses the latter's explicit runtime provenance and the former's
unchanged observations. The complete native path is
\begin{quote}\small
\path{verification/evidence/five-contracts-computer10-20261104/runtime-final/towers/check/simulation.csv}.
\end{quote}
Its SHA-256 is
\begin{quote}\small
\path{18044600a574e19d0bdd2576fdba7782c8b43d80cd6b3bae9a30fba2b1ea35b5}.
\end{quote}
The adjacent \path{runtime.json} pins source hashes, executable, commands,
Ghost policy, exit statuses and validation logs. The parallel
\path{towers/ignore/} directory records the Ghost Ignore execution. These
are replays of the same stream, not additional Monte Carlo replicates.

\section{Generator, transformation and exact counting}

The native generator updates its integer state by the Park--Miller recurrence
\[
 s_{t+1}=48271s_t\pmod{2147483647},\qquad
 u_{t+1}=s_{t+1}/2147483647.
\]
Two successive values \(u,v\) feed the Box--Muller transformation
\[
 r=\sqrt{-2\log u},\quad \theta=2\pi v,\quad
 X=r\cos\theta,\quad E=r\sin\theta,
 \qquad Y=0.7X+\sqrt{0.51}\,E.
\]
The finite-state recurrence, conversion and checked normal-value bounds are
implementation properties. The exact continuous-normal interpretation needs
the additional idealization that \(u,v\) are independent continuous uniforms;
SPARK does not prove that idealization for deterministic states.

\begin{proposition}[Ideal Box--Muller law]\label{tw:box-muller}
If \(U,V\) are independent uniform variables on \((0,1)\), the displayed
transform produces independent standard normals \(X,E\).
\end{proposition}
\begin{proof}
Set \(R=\sqrt{-2\log U}\) and \(\Theta=2\pi V\).
For \(r>0\),
\(P(R\leq r)=P(U\geq e^{-r^2/2})=1-e^{-r^2/2}\), so the radial
density is \(re^{-r^2/2}\). The angle has density \(1/(2\pi)\), and
the two variables remain independent because they are measurable functions
of separate independent inputs. Thus their joint density is
\(re^{-r^2/2}/(2\pi)\). The polar-to-Cartesian map has absolute Jacobian
\(r\). Dividing by this Jacobian gives, away from the origin and the
angular seam (sets of probability zero),
\[
 f_{X,E}(x,e)=\frac{1}{2\pi}e^{-(x^2+e^2)/2}
             =\phi(x)\phi(e).
\]
The change-of-variables theorem and density factorization therefore give
standard-normal margins and independence. These are the density and
independence principles used for Gaussian transformations in
\citet[Chapter 3]{gallager2013}; \cref{tw:gaussian} then proves that the
last affine transformation gives correlation \(0.7\).
\end{proof}

The SPARK driver computes the four pairs of cutoffs before resetting the
generator. For each generated pair it calls the same strict-threshold
classifier under each rule, increments exactly one cell and maintains total
count equal to the number of observations processed. After the loop it
extracts the marginal and union counts, joint fraction and two conditional
fractions. A zero conditional denominator is represented as missing in the
CSV, rather than as a zero probability. The four-cell partition and all these
integer identities are proved in \cref{tw:simulation}; their validity does
not require the normal model or random independence.

\section{All native event cells and their interpretation}

The following are the complete native cell observations, transcribed from
the pinned CSV. Every row sums to one million; equal marginal cutoffs on a
given test reuse exactly the same passing indicators across rules.
\begin{center}\small
\begin{tabular}{lrrrr}
\toprule
Rule & Neither \(n_{00}\) & Only \(A\), \(n_{10}\) & Only \(B\), \(n_{01}\) & Both \(n_{11}\)\\
\midrule
E2    & 966622 & 13665 & 13474 & 6239\\
E0.1  & 998193 & 825   & 815   & 167\\
Rare  & 999936 & 28    & 32    & 4\\
Mixed & 979754 & 19264 & 342   & 640\\
\bottomrule
\end{tabular}
\end{center}
The Dask workflow's retained-Ada table derives the denominators and
forward conditional rates from these same cells:

\begingroup\small
\setlength{\tabcolsep}{4pt}
\Needspace{6\baselineskip}
\noindent\textbf{Retained Ada run: one million pairs under each rule}
\begin{longtable}{rrrrrr}
\toprule
Rule & $N_A$ & $N_B$ & Both & Union & $P(B|A)$ \\
\midrule
\endfirsthead
Rule & $N_A$ & $N_B$ & Both & Union & $P(B|A)$ \\
\midrule
\endhead
\bottomrule
\endfoot
E2 & $19904$ & $19713$ & $6239$ & $33378$ & $0.313455$ \\
E0.1 & $992$ & $982$ & $167$ & $1807$ & $0.168347$ \\
Rare & $32$ & $36$ & $4$ & $64$ & $0.125000$ \\
Mixed & $19904$ & $982$ & $640$ & $20246$ & $0.032154$ \\
\end{longtable}
\endgroup


For E2, \(33378/10^6=0.033378\) pass at least one test, while
\(6239/10^6=0.006239\) pass both. The two repeat-qualification rates are
\(6239/19904\approx0.313455\) and
\(6239/19713\approx0.316492\); their finite-sample difference is compatible
with different observed marginal counts even though the model is symmetric.
For E0.1, the union is \(1807/10^6=0.001807\), and the corresponding
conditional rates are \(167/992\approx0.168347\) and
\(167/982\approx0.170061\). These are different questions from its joint
frequency \(0.000167\).

For Rare, only four pairs pass both. Their frequency is
\(4\times10^{-6}\), against the Gaussian reference
\(j\approx2.3125\times10^{-6}\). The reported 95\% Wilson interval is
approximately \([1.55552\times10^{-6},1.02859\times10^{-5}]\), a wide
range relative to the target. The conditional rates \(4/32=0.125\) and
\(4/36\approx0.111111\) depend on only 32 and 36 selected passers.
Neither four successes nor those small denominators support fine decimal
precision about population repeat qualification. Under the ideal model,
the expected native joint count is only \(10^6j\approx2.3125\).
The binomial moments and relative-error calculation are proved in
\cref{sm:binomial,sm:rare}; the Wilson inversion and conditional-binomial
law are proved in \cref{sm:wilson,sm:conditional-count}.

For Mixed, the native union is \(20246/10^6=0.020246\).
Its conditional rates \(640/19904\approx0.032154\) and
\(640/982\approx0.651731\) have the same numerator and sharply different
denominators. The model targets are approximately \(0.0329294\) and
\(0.658589\). This reproduces the direction and scale of the denominator
effect without turning the finite estimates into analytic equalities.
Every displayed rate is obtained by the event identities in
\cref{tw:events,tw:simulation}, not by reading overlap from a graph.

\section{Direct comparison with the Dask experiment}

For the four shared rules, the native denominator is one million and the
Dask pooled denominator is four million. The directly comparable joint and
union frequencies are:
\begin{center}\small
\begin{tabular}{lrrrr}
\toprule
 & \multicolumn{2}{c}{SPARK, \(T_S=10^6\)}
 & \multicolumn{2}{c}{Dask, \(T_D=4\times10^6\)}\\
Rule & Joint & Union & Joint & Union\\
\midrule
E2    & 0.006239 & 0.033378 & 0.006247   & 0.033634\\
E0.1  & 0.000167 & 0.001807 & 0.000157   & 0.00180125\\
Rare  & 0.000004 & 0.000064 & 0.00000275 & 0.00006150\\
Mixed & 0.000640 & 0.020246 & 0.00065675 & 0.02034325\\
\bottomrule
\end{tabular}
\end{center}
These are descriptive cross-implementation comparisons, not tests that the
two finite streams are independent. The Dask four-million Rare row has 11
joint successes and an exact binomial interval approximately
\([1.37279\times10^{-6},4.92050\times10^{-6}]\); it improves event
information over four successes while retaining substantial relative
uncertainty. The native Wilson interval and this exact interval use different
constructions, whose proofs and coverage distinctions appear in
\cref{sm:wilson,sm:clopper-pearson}. For any \emph{one} fixed event with
ideal IID pairs, its frequency variance is \(p(1-p)/T\) by
\cref{sm:binomial}; replacing \(T\) by \(4T\) halves its reference standard
error. This sampling-law comparison does not assert exact deterministic
generator quality or an independent-stream standard error for the difference.

Dask also simulates OR2, which the native campaign only calibrates
numerically. Its pooled observed union is \(0.019815\), against the target
\(0.02\). Thus the report retains both the native calculation of the
calibrated cutoff and Dask's finite-sample check at that fixed cutoff.
The agreement of the two implementations concerns model-level targets;
neither matching decimal values nor different seeds establishes real-test
Gaussianity, empirical repeatability or psychometric validity.

\section{What the native verification establishes}

The 4 November Towers portfolio records 696 discharged obligations:
230 flow and 466 prover checks, zero unproved and zero justified checks.
It reports 89 entities with native \texttt{all} status and 16 with
\texttt{spec} status. These are dated, source-pinned facts in
\path{verification/evidence/five-contracts-computer10-20261104/towers/portfolio/result.json}.
The proved obligations cover the bounded counter core, classification,
recurrence and conversion, driver loops, result assembly and specified
numerical arithmetic. The source inventory records 99 project subprograms
and 557 explicit contract entries, with a minimum of five. These counts
describe that campaign and do not establish adequacy or independence of
all its contracts.

Checked completions for elementary functions, foreign normal-tail evaluation,
numerical acceptance guards and physical I/O remain explicit trusted
boundaries. Proved callers and normal-return receipts do not prove backend
totality, Gaussian randomness, integration error, correct file persistence
or validity of the statistical model. The retained native checks are not
claimed as a new proof of subsequently refactored library source, or as a
proof of the Python/Dask implementation.

The retained runtime records separately show Ghost Check and Ghost Ignore
builds with ordinary assertions enabled, byte-identical CSVs, successful
validation in normal and optimized Python, and exact replay of all 16
event-cell counts by the separate Python implementation. Numerical outputs
agree with the independent SciPy reference within the recorded absolute
tolerance \(10^{-10}\). Runtime bookkeeping, replay agreement, approximate
numerical validation and static SPARK proof therefore supply different
evidence. None is relabeled as a proof of finite-sample independence or of
extreme-tail accuracy for actual people.
